\chapter{Dihydropyrimidine Dehydrogenase (DPD/DPYD) Testing}

\section*{What Is This Test?}
DPD testing identifies reduced activity of dihydropyrimidine dehydrogenase, the main enzyme that breaks down fluoropyrimidine medicines. Testing may use DPYD genotyping, a direct enzyme-activity assay, or both.

\section*{Alternative Names}
\begin{itemize}
  \item DPD Deficiency Test
  \item DPYD Genotyping
  \item Dihydropyrimidine Dehydrogenase Deficiency
  \item Fluoropyrimidine Toxicity Risk Test
\end{itemize}
\section*{Why This Test Is Ordered}
Testing is used before starting fluoropyrimidines such as 5-fluorouracil, capecitabine, or tegafur where recommended by local guidance. Reduced DPD activity increases the risk of severe or fatal toxicity, including mucositis, diarrhea, neutropenia, neurotoxicity, and cardiotoxicity.

\section*{How This Test Is Performed}
Genotyping uses a blood or saliva sample to detect clinically important DPYD variants. Phenotype testing measures DPD activity or uracil-related metabolites in blood and may identify risk not captured by a limited genetic panel.

\section*{Preparation, Risks, and Limitations}
No fasting is required. Risks are limited to sample collection. A normal limited genotype panel does not exclude every cause of reduced enzyme activity, and results must be interpreted with ancestry, assay coverage, the planned medicine, dose, and specialist pharmacogenetic guidance.

\section*{Understanding Results}
Normal, intermediate, or poor-metabolizer results may lead to standard dosing, a reduced starting dose, an alternative medicine, or enhanced monitoring according to current clinical guidance. Testing does not predict every adverse effect and does not replace clinical monitoring.

\section*{References}
Clinical Pharmacogenetics Implementation Consortium guidance on DPYD genotype and fluoropyrimidine dosing.

\clearpage
\section*{Understanding Results and Follow-up}
The laboratory report should identify the specimen, method, units, reference interval or decision limit, and whether the result is positive, negative, abnormal, or inconclusive. Reference intervals vary with age, sex, pregnancy, specimen type, assay, and laboratory; a result outside a range is not automatically a diagnosis, and a result within a range does not always exclude disease. Discuss unexpected results with the ordering clinician, who can decide whether repeat or confirmatory testing, treatment, or referral is appropriate. Do not change medicines or delay urgent care based on an isolated result.


\section*{Regulatory Status and Authorization Date}
This chapter describes a test category, not a specific commercial product. FDA, CDSCO, EU IVDR, and MHRA approval or registration dates therefore cannot be assigned without the assay manufacturer, product name, instrument, and intended use. For laboratory-developed tests, an individual agency approval date may not exist. See the Preface for the interpretation of regulatory status.

\clearpage
